\documentclass[../../main.tex]{subfiles}
\begin{document}

{\bf[解]}

\subparagraph{(1)}

与えられた関数
\begin{align}
  f(x) = x^4+x^3+\dfrac{1}{2}x^2+\dfrac{1}{6}x+\dfrac{1}{24} \label{eq:1}
\end{align}
について考える．
一階微分，二階微分を計算すると
\begin{align}
  f'(x)  &=4x^3+3x^2+x+\dfrac{1}{6} \label{eq:2}\\
  f''(x) &=12(x+\dfrac{1}{4})^2+\dfrac{1}{4}>0
\end{align}
である．$f''(x)>0$より$f'(x)$は単調増加である．
また，$f'(x)$は$x\to\pm\infty$の時に$f'(x)\to\pm\infty$だから，
$f'(t)=0$となる$t$がただ一つ存在する．
これを$x=\beta$とすると，$f(x)$の増減表は下表になる．

\begin{table}[htb]
  \centering
  \caption{$f(x)$ の増減表}
  \renewcommand{\arraystretch}{1.4}
  \begin{tabular}{c|ccccc}
    $x$     & $\dots$  & $\beta$         & $\dots$  \\ \hline
    $f'(x)$ & $-$      & $0$             & $+$      \\ \hline
    $f(x)$  & $\searrow$ & \text{極小}   & $\nearrow$
  \end{tabular}
\end{table}

よって$f(x) \ge f(\beta)$だから$f(\beta)>0$を示せば良い．
\cref{eq:1,eq:2}より
\begin{align}
  f(x) = (\dfrac{1}{4}x+\dfrac{1}{16})f'(x)+\dfrac{1}{32}(2x^2+2x+1)
\end{align}
だから$x=\beta$を代入して$f'(\beta)=0$より
\begin{align}
  f(\beta)
  &=\dfrac{1}{32}(2\beta^2+2\beta+1) \\
  &=\dfrac{1}{32}\left\{2\left(\beta+\dfrac{1}{2}\right)^2+\dfrac{1}{2}\right\} >0
\end{align}
となって$f(\beta)>0$だから題意は示された．$\Box$

\subparagraph{(2)}

与えられた関数
\begin{align}
  g(x) = x^5+x^4+\dfrac{1}{2}x^3+\dfrac{1}{6}x^2+\dfrac{1}{24}x+\dfrac{1}{120} \label{eq:3}
\end{align}
について考える．
一階微分，二階微分を計算すると
\begin{align}
  g'(x)&=5x^4+4x^3+\dfrac{3}{2}x^2+\dfrac{1}{3}x+\dfrac{1}{24} \\
  g''(x)&=20x^3+12x^2+3x+\dfrac{1}{3}  \\
  g'''(x)&=60\left(x+\dfrac{1}{5}\right)^2+\dfrac{3}{5}>0
\end{align}
である．
従って前問と同様に$g''(\gamma)=0$なる実数$\gamma$がただひとつ存在し，
$g'(x)$の増減表は以下のようになる．

\begin{table}[htb]
  \centering
  \caption{$g'(x)$ の増減表}
  \label{tab:fx_monotonicity}
  \renewcommand{\arraystretch}{1.4}
  \begin{tabular}{c|ccccc}
    $x$     & $\dots$     & $\gamma$        & $\dots$    \\ \hline
    $g''(x)$ & $-$        & $0$             & $+$        \\ \hline
    $g'(x)$  & $\searrow$ & \text{極小}     & $\nearrow$
  \end{tabular}
\end{table}

よって
\begin{align}
  g'(x)
  &\ge g'(\gamma) \\
  &=\left(\dfrac{1}{4}\gamma+\dfrac{1}{20}\right)g''(\gamma)+40\left\{6\left(\gamma+\dfrac{1}{3}\right)^2+\dfrac{1}{3}\right\} \\
  &=40\left\{6\left(\gamma+\dfrac{1}{3}\right)^2+\dfrac{1}{3}\right\}>0
\end{align}
である．
従って$g'(x)>0$で$g(x)$は単調増加である．
これと$g(0)>0$，$g(-1)=\dfrac{-11}{30}$から
$g(x)=0$はただ一つの実数解を$-1<x<0$に持ち，題意は示された．$\Box$

{\bf[別解]}

\subparagraph{(1)}

$f(x)$を変形すると
\begin{align}
  f(x)=\left(x^2+\dfrac{1}{2}x\right)^2+\dfrac{1}{4}\left(x+\dfrac{1}{3}\right)^2+\dfrac{1}{72}
\end{align}
より，確かに$f(x)>0$である．よって示された．$\Box$

\subparagraph{(2)}
前問と同様に
\begin{align}
  g'(x)
  &=5x^4+4x^3+\dfrac{3}{2}x^2+\dfrac{1}{3}x+\dfrac{1}{24}     \\
  &=5\left(x^2+\dfrac{2}{5}x\right)^2+\dfrac{7}{10}\left(x+\dfrac{5}{21}\right)^2+\dfrac{3}{1512}
\end{align}
ゆえ$g'(x)>0$である．これと$g(-1)<0$および$g(0)>0$から示すべきことは明らか．$\Box$

% TODO: sin/cosのTaylor展開の解説

\newpage
\end{document}
